% ============================================================
% Chapter 14 — Reporting Methodology
% ============================================================
\chapter{Reporting Methodology}
\label{ch:reporting}

\noteblock{A penetration test is only as valuable as its report. A well-written report translates technical findings into business risk, enabling developers and management to understand and prioritize remediation.}

% ─────────────────────────────────────────────────────────
\section{Report Structure}

A professional penetration test report consists of two parts:

\begin{enumerate}
  \item \textbf{Executive Summary} — for management: business risk, severity overview, key recommendations
  \item \textbf{Technical Findings} — for developers: vulnerability details, evidence, remediation steps
\end{enumerate}

% ─────────────────────────────────────────────────────────
\section{Severity Classification}

Always use a consistent severity classification system. CVSS v3.1 is the industry standard.

\begin{table}[H]
\centering
\caption{Severity classification}
\begin{tabularx}{\textwidth}{>{\bfseries}lXll}
\toprule
Severity & Description & CVSS Range & Remediate Within \\
\midrule
Critical & Immediate full compromise possible & 9.0--10.0 & 24--48 hours \\
High & Significant risk, likely to be exploited & 7.0--8.9 & 1 week \\
Medium & Exploitable with some conditions & 4.0--6.9 & 1 month \\
Low & Limited impact, low exploitability & 0.1--3.9 & 3 months \\
Informational & No direct risk, best practice & N/A & Next review cycle \\
\bottomrule
\end{tabularx}
\end{table}

% ─────────────────────────────────────────────────────────
\section{Finding Report Format}

Every finding must include the following fields:

\findingblock{
\begin{tabular}{ll}
\textbf{Finding Title:} & [Short, descriptive name] \\
\textbf{Severity:} & Critical / High / Medium / Low / Informational \\
\textbf{CVSS v3.1 Score:} & [X.X] \\
\textbf{CVSS Vector:} & [AV:N/AC:L/PR:N/UI:N/S:U/C:H/I:H/A:N] \\
\textbf{CWE:} & CWE-XXX \\
\textbf{OWASP:} & A0X:2021 \\
\textbf{Affected Component:} & [URL, endpoint, parameter] \\
\textbf{Discovered:} & [Date] \\
\end{tabular}\\[8pt]

\textbf{Description:}\\
[2--4 sentences explaining the vulnerability in plain English. Assume the reader is a developer, not a security expert.]\\[6pt]

\textbf{Steps to Reproduce:}\\
1. [Step 1]\\
2. [Step 2]\\
3. [Step 3]\\[6pt]

\textbf{Evidence:}\\
[Include the exact HTTP request and response from Burp Suite]\\
\begin{lstlisting}[style=http]
REQUEST:
GET /vulnerable-endpoint?param=payload HTTP/1.1
Host: target.com
Authorization: Bearer <token>

RESPONSE:
HTTP/1.1 200 OK
[Response showing vulnerability evidence]
\end{lstlisting}

\textbf{Impact:}\\
[Business impact — what can an attacker achieve? Be specific: data exfiltration, account takeover, financial loss.]\\[6pt]

\textbf{Recommendation:}\\
[Actionable, specific remediation steps. Reference secure coding patterns.]\\[6pt]

\textbf{References:}\\
{[}OWASP link{]}, {[}CWE link{]}, {[}CVE if applicable{]}
}

% ─────────────────────────────────────────────────────────
\section{Executive Summary Template}

\begin{tcolorbox}[findingstyle, title={Executive Summary Template}]
\textbf{Assessment Overview}\\[4pt]
A web application security assessment was conducted against \texttt{[target]} between \texttt{[start date]} and \texttt{[end date]}. The assessment identified \texttt{[N]} vulnerabilities across \texttt{[N]} endpoints.\\[6pt]

\textbf{Risk Summary}\\[4pt]
\begin{tabular}{ll}
Critical: & [N] \\
High: & [N] \\
Medium: & [N] \\
Low: & [N] \\
Informational: & [N] \\
\textbf{Overall Risk:} & \textbf{[Critical/High/Medium/Low]} \\
\end{tabular}\\[6pt]

\textbf{Key Findings}\\[4pt]
The most significant findings were:
\begin{itemize}
  \item [Finding 1 in one sentence]
  \item [Finding 2 in one sentence]
  \item [Finding 3 in one sentence]
\end{itemize}

\textbf{Immediate Actions Required}\\[4pt]
The following must be addressed before the next production deployment:
\begin{itemize}
  \item [Critical action 1]
  \item [Critical action 2]
\end{itemize}
\end{tcolorbox}

% ─────────────────────────────────────────────────────────
\section{Evidence Collection Best Practices}

\begin{itemize}
  \item \textbf{Always capture the full HTTP request and response} — partial evidence is insufficient
  \item \textbf{Use Burp's copy-as-curl} for easily reproducible PoC commands
  \item \textbf{Include timestamps} — important for proving the finding occurred at a specific time
  \item \textbf{Never alter real data} — use test accounts and synthetic payloads only
  \item \textbf{Document the testing account} — which account was used to find each issue
  \item \textbf{Save Burp project files} (Pro) or export HTTP logs (Community) as backup evidence
\end{itemize}

% ─────────────────────────────────────────────────────────
\section{Remediation Verification}

After developers apply fixes, the tester should retest each finding:

\begin{table}[H]
\centering
\caption{Remediation verification outcomes}
\begin{tabularx}{\textwidth}{>{\bfseries}lX}
\toprule
Outcome & Action \\
\midrule
Fixed & Mark as \textbf{Resolved} — document the retest date and result \\
Partially fixed & Mark as \textbf{Partially Remediated} — document what remains \\
Not fixed & Mark as \textbf{Open} — escalate if critical \\
Introduced new issue & Document as \textbf{New Finding} — log separately \\
\bottomrule
\end{tabularx}
\end{table}

\warnblock{Retesting should use the same methodology as the original test. A finding should only be marked resolved if the tester has confirmed the exact original exploit no longer works — not just because the developer says it is fixed.}
